\section{Discussion}

MetaShift-Bench makes a weakly labelled question auditable rather than
pretending that it is fully identified. The critical design choice is to retain
metadata anchors, availability constraints, donor rules, algorithm comparisons,
and failures in the same evidence chain. A result from a well-matched target is
useful only when its selection requirements, reference composition, uncertainty
scope, and external evidence limitations remain visible.

The synthetic experiment is the source of known truth in this project. It
shows that cross-site counterfactual construction can separate some local and
matched regional patterns better than single-series baselines, while also
showing that MetaShift does not achieve a stable aggregate advantage over
standard synthetic control. Reporting that negative result prevents the
benchmark from being reinterpreted as a method-promotion exercise.

The real-event audit has a complementary role. It asks which metadata anchors
can be evaluated under transparent, predeclared inputs and which cannot. A
complete residual comparison is not evidence of a physical measurement error.
Conversely, a donor-insufficient event is not evidence that no change occurred;
it is an abstention because the protocol cannot build a qualified geographic
reference. This distinction makes the full ledger more informative than a
curated set of striking examples.

The evidence tiers are best understood as a prioritization tool for future
human review. A supported candidate has a coherent pattern across the
predefined diagnostics, but contextual review, station history, and direct
technical records would still be needed before assigning a physical
explanation. The same caution applies to POC comparisons and public document
searches: they characterize available context rather than validate an
instrument-change hypothesis.
